\documentclass[12pt,a4paper]{article}

\usepackage[a4paper,top=25mm,bottom=25mm,left=22mm,right=22mm]{geometry}
\usepackage{fontspec}
\usepackage{xepersian}
\usepackage{longtable}
\usepackage{booktabs}
\usepackage{array}
\usepackage{xcolor}
\usepackage{hyperref}
\usepackage{enumitem}
\usepackage{fancyhdr}
\usepackage{titlesec}
\usepackage{fancyvrb}

\IfFontExistsTF{Vazirmatn}{
  \settextfont{Vazirmatn}
}{
  \IfFontExistsTF{Tahoma}{
    \settextfont{Tahoma}
  }{
    \settextfont{Arial}
  }
}
\IfFontExistsTF{Consolas}{
  \setlatintextfont{Consolas}
}{
  \setlatintextfont{Courier New}
}

\definecolor{primary}{HTML}{1F4E79}
\definecolor{softgray}{HTML}{F3F6F8}
\definecolor{warn}{HTML}{A15C00}

\hypersetup{
  colorlinks=true,
  linkcolor=primary,
  urlcolor=primary,
  citecolor=primary,
  pdftitle={مستند DevOps سامانه AI Chat SaaS},
  pdfauthor={Mobin Yousefi}
}

\pagestyle{fancy}
\fancyhf{}
\rhead{مستند DevOps}
\lhead{AI Chat SaaS}
\cfoot{\thepage}
\renewcommand{\headrulewidth}{0.4pt}

\titleformat{\section}{\Large\bfseries\color{primary}}{\thesection}{0.7em}{}
\titleformat{\subsection}{\large\bfseries\color{primary}}{\thesubsection}{0.7em}{}
\titleformat{\subsubsection}{\normalsize\bfseries\color{primary}}{\thesubsubsection}{0.7em}{}
\setlist[itemize]{topsep=4pt,itemsep=2pt,parsep=0pt}
\setlist[enumerate]{topsep=4pt,itemsep=2pt,parsep=0pt}
\renewcommand{\arraystretch}{1.35}
\setlength{\LTpre}{6pt}
\setlength{\LTpost}{6pt}
\setlength{\parindent}{0pt}
\setlength{\parskip}{6pt}
\setlength{\headheight}{15pt}

\newcommand{\en}[1]{\lr{#1}}
\newcommand{\code}[1]{\lr{\texttt{\detokenize{#1}}}}
\newcolumntype{R}[1]{>{\raggedleft\arraybackslash}p{#1}}
\newenvironment{codeblock}{\VerbatimEnvironment\begin{latin}\begin{Verbatim}[fontsize=\small,frame=single,rulecolor=\color{softgray}]}{\end{Verbatim}\end{latin}}

\begin{document}

\begin{titlepage}
\centering
\vspace*{22mm}
{\Huge\bfseries\color{primary} مستند DevOps سامانه AI Chat SaaS\par}
\vspace{10mm}
{\Large بخش ۱۲ راهنمای مستندات پروژه\par}
\vspace{20mm}
{\large \textbf{نسخه سند:} 1.0\par}
\vspace{3mm}
{\large \textbf{تاریخ تدوین:} ۱۵ جولای ۲۰۲۶\par}
\vspace{3mm}
{\large \textbf{تهیه‌کننده:} مبین یوسفی\par}
\vspace{3mm}
{\large \textbf{مسیر خروجی:} \code{doc/LaTeX/12-DevOps/12-DevOps.tex}\par}
\vspace{3mm}
{\large \textbf{منابع بررسی‌شده:} \code{Rahnama.md}، \code{ArchitectureForMobin.md}، \code{ReportForSale.md}، \code{README.md} و کدهای پروژه\par}
\vfill
{\large مستند راهبری توسعه، build، deployment، rollback، monitoring، logging و آماده‌سازی production}
\end{titlepage}

\tableofcontents
\newpage

\section{هدف و دامنه سند}
این سند فرآیندهای DevOps سامانه \en{AI Chat SaaS} را برای توسعه، build، release، deployment، rollback، monitoring و نگهداری عملیاتی توضیح می‌دهد. سند بر اساس وضعیت واقعی repository نوشته شده است: backend با PHP خام، frontend با Next.js، widget به‌صورت JavaScript قابل انتشار، پایگاه داده MySQL و تنظیمات محیطی مبتنی بر \code{.env}.

در وضعیت فعلی، فایل رسمی برای Docker، Docker Compose، Kubernetes یا CI/CD در repository وجود ندارد. بنابراین این سند دو بخش را هم‌زمان پوشش می‌دهد: وضعیت موجود و الگوی پیشنهادی برای production. هرجا قابلیت هنوز پیاده‌سازی نشده، صریحا با عنوان «پیشنهادی» یا «نیازمند تکمیل» مشخص شده است.

\section{نمای کلی اجزای قابل استقرار}
\begin{longtable}{R{34mm} R{60mm} R{58mm}}
\toprule
\textbf{جزء} & \textbf{مسیر/فناوری} & \textbf{مسئولیت عملیاتی} \\
\midrule
\endhead
Backend API & \code{backend/}، PHP، PDO، MySQL & سرو APIها، احراز هویت، ویجت، پنل، AI، upload و دسترسی دیتابیس. \\
Frontend Panel & \code{frontend/}، Next.js، React & پنل super admin، customer admin و agent؛ build و اجرای Node/Next. \\
Widget & \code{widget/dist/widget.js} & فایل JavaScript قابل embed روی سایت مشتری؛ باید از دامنه یا CDN پایدار سرو شود. \\
Database & MySQL با migrationهای SQL & نگهداری tenantها، کاربران، سایت‌ها، گفتگوها، دانش AI، rate limit و audit. \\
Uploads & \code{backend/uploads} یا \code{UPLOAD_STORAGE_ROOT} & نگهداری پیوست‌ها و تصاویر؛ نیازمند backup و hardening وب‌سرور. \\
Test Shop & \code{test-shop/} و \code{widget/test.html} & محیط تست/دمو برای نصب ویجت و سناریوهای نمایشی. \\
\bottomrule
\end{longtable}

\section{وضعیت فعلی Repository}
\subsection{فایل‌ها و ابزارهای موجود}
\begin{itemize}
  \item \code{frontend/package.json}: شامل اسکریپت‌های \code{dev}، \code{build}، \code{start} و \code{lint}.
  \item \code{frontend/next.config.ts}: هدرهای امنیتی frontend را تعریف می‌کند.
  \item \code{backend/.env.example}: نمونه متغیرهای محیطی backend، دیتابیس، JWT، upload و CORS.
  \item \code{backend/.htaccess}: جلوگیری از directory listing، محافظت از فایل‌های env/sql/log/backup و خاموش کردن display errors در mod\_php.
  \item \code{backend/database/migrations/2026_07_07_ai_knowledge_tables.sql}: migration مربوط به جدول‌های دانش AI.
  \item \code{backend/api/test-db.php}: health check ساده برای اتصال دیتابیس و شمارش جدول \code{plans}.
  \item \code{widget/src/widget.js} و \code{widget/dist/widget.js}: سورس و خروجی قابل انتشار ویجت.
\end{itemize}

\subsection{فایل‌ها و فرآیندهای غایب}
\begin{itemize}
  \item \code{Dockerfile} و \code{docker-compose.yml} در repository وجود ندارد.
  \item pipeline رسمی CI/CD مانند GitHub Actions، GitLab CI یا Jenkinsfile وجود ندارد.
  \item migration runner خودکار دیده نشد؛ migrationها باید فعلا دستی یا با اسکریپت جدا اجرا شوند.
  \item health check جامع برای همه سرویس‌ها وجود ندارد؛ فقط \code{test-db.php} برای دیتابیس موجود است.
  \item monitoring، tracing و centralized logging در repo پیاده‌سازی نشده‌اند.
\end{itemize}

\section{محیط‌های عملیاتی}
\begin{longtable}{R{30mm} R{58mm} R{66mm}}
\toprule
\textbf{محیط} & \textbf{هدف} & \textbf{قواعد پیشنهادی} \\
\midrule
\endhead
Local & توسعه روی ماشین برنامه‌نویس & \code{APP_ENV=local}، \code{APP_DEBUG=true}، دیتابیس محلی، CORS برای localhost. \\
Staging & تست قبل از انتشار & داده تست، دامنه جدا، \code{APP_DEBUG=false}، اجرای migration و smoke test. \\
Production & سرویس واقعی مشتریان & HTTPS اجباری، رازهای قوی، backup، monitoring، محدودیت CORS و rollback آماده. \\
\bottomrule
\end{longtable}

\subsection{متغیرهای محیطی backend}
متغیرهای اصلی از \code{backend/.env} خوانده می‌شوند. فایل \code{.env} نباید commit شود و فقط \code{.env.example} باید در repository باقی بماند.

\begin{codeblock}
APP_ENV=production
APP_DEBUG=false
APP_URL=https://app.example.com
FRONTEND_URL=https://app.example.com
API_URL=https://api.example.com/backend/api
WIDGET_SCRIPT_URL=https://cdn.example.com/widget.js

DB_HOST=127.0.0.1
DB_NAME=ai_chat_saas
DB_USER=ai_chat_saas_user
DB_PASS=strong-password

JWT_SECRET=long-random-secret-at-least-32-characters
PANEL_ALLOWED_ORIGINS=https://app.example.com
WIDGET_ALLOWED_ORIGINS=https://customer-one.com
WIDGET_ALLOW_EMPTY_ORIGIN=false

UPLOAD_STORAGE_ROOT=/var/www/ai-chat-saas-storage
UPLOAD_PUBLIC_URL=https://api.example.com
\end{codeblock}

\subsection{متغیرهای frontend}
frontend از \code{NEXT_PUBLIC_API_BASE_URL} استفاده می‌کند. اگر این مقدار تنظیم نشود، مقدار local زیر استفاده می‌شود:
\begin{codeblock}
http://localhost/ai-chat-saas/backend/api
\end{codeblock}

برای production:
\begin{codeblock}
NEXT_PUBLIC_API_BASE_URL=https://api.example.com/backend/api
\end{codeblock}

\section{اجرای محلی}
\subsection{پیش‌نیازها}
\begin{itemize}
  \item PHP با افزونه‌های PDO و MySQL.
  \item MySQL یا MariaDB.
  \item Node.js و npm برای frontend.
  \item وب‌سرور محلی مانند Apache/XAMPP/WAMP یا PHP built-in server برای backend.
\end{itemize}

\subsection{راه‌اندازی backend}
\begin{enumerate}
  \item فایل \code{backend/.env.example} به \code{backend/.env} کپی و مقادیر دیتابیس تنظیم شود.
  \item دیتابیس \code{ai_chat_saas} ساخته شود.
  \item schema پایه و migrationهای موجود اجرا شوند. در repo فعلی migration دانش AI در \code{backend/database/migrations} قرار دارد.
  \item API از طریق وب‌سرور محلی در مسیری مانند \code{http://localhost/ai-chat-saas/backend/api} سرو شود.
  \item اتصال دیتابیس با endpoint زیر تست شود:
\end{enumerate}

\begin{codeblock}
GET /backend/api/test-db.php
\end{codeblock}

\subsection{راه‌اندازی frontend}
\begin{codeblock}
cd frontend
npm install
npm run dev
\end{codeblock}

برای build production:
\begin{codeblock}
cd frontend
npm ci
npm run build
npm run start
\end{codeblock}

\subsection{راه‌اندازی widget}
ویجت فعلی ابزار build جداگانه ندارد و فایل قابل انتشار در \code{widget/dist/widget.js} موجود است. در محیط توسعه می‌توان از \code{widget/test.html} یا embed زیر استفاده کرد:

\begin{codeblock}
<script
  src="https://cdn.example.com/widget.js"
  data-site-key="SITE_PUBLIC_KEY"
  data-api-base="https://api.example.com/backend/api">
</script>
\end{codeblock}

\section{Git و Branch Strategy}
\subsection{قواعد پایه Git}
\begin{itemize}
  \item شاخه اصلی پایدار باید \code{main} یا \code{master} باشد.
  \item تغییرات مستقیم روی شاخه production ممنوع باشد.
  \item هر قابلیت یا اصلاح با branch جدا و merge request انجام شود.
  \item فایل‌های \code{.env}، dump دیتابیس، uploadها، logها و buildهای موقت commit نشوند.
  \item پیام commit باید کوتاه، دقیق و قابل ردیابی باشد.
\end{itemize}

\subsection{Git Flow پیشنهادی}
\begin{longtable}{R{34mm} R{100mm}}
\toprule
\textbf{Branch} & \textbf{کاربرد} \\
\midrule
\endhead
\code{main} & نسخه production یا آماده release. فقط از طریق merge request و پس از تایید. \\
\code{develop} & تجمیع تغییرات آماده staging. اگر تیم کوچک است، می‌تواند حذف شود و فقط feature branch به main merge شود. \\
\code{feature/*} & توسعه قابلیت جدید مانند \code{feature/ai-crawl-monitoring}. \\
\code{fix/*} & رفع bugهای عادی. \\
\code{hotfix/*} & اصلاح فوری production با مسیر کوتاه‌تر و release سریع. \\
\code{release/*} & آماده‌سازی نسخه، freeze تغییرات، اجرای تست و اصلاحات پایانی. \\
\bottomrule
\end{longtable}

\subsection{Merge Request}
هر merge request باید حداقل شامل موارد زیر باشد:
\begin{itemize}
  \item شرح تغییر و دلیل آن.
  \item فهرست endpointها یا صفحات تحت تاثیر.
  \item وضعیت migration دیتابیس.
  \item نتیجه build/lint/test.
  \item ریسک rollback.
  \item تایید اینکه secret یا داده حساس commit نشده است.
\end{itemize}

\section{Versioning و Release}
برای محصول SaaS، نسخه‌بندی پیشنهادی \en{Semantic Versioning} است:
\begin{codeblock}
MAJOR.MINOR.PATCH
1.4.2
\end{codeblock}

\begin{itemize}
  \item \code{MAJOR}: تغییر ناسازگار در API، دیتابیس یا قرارداد widget.
  \item \code{MINOR}: قابلیت جدید سازگار با نسخه قبل.
  \item \code{PATCH}: bug fix، hardening یا تغییر بدون اثر روی قرارداد عمومی.
\end{itemize}

هر release باید tag داشته باشد:
\begin{codeblock}
git tag -a v1.0.0 -m "Release v1.0.0"
git push origin v1.0.0
\end{codeblock}

\section{Build و Artifact}
\subsection{Backend}
backend فعلی build ندارد، اما artifact آن باید شامل فایل‌های PHP، \code{backend/api}، \code{backend/includes}، \code{backend/config}، \code{backend/public}، migrationها و فایل \code{.htaccess} باشد. فایل‌های زیر نباید وارد artifact شوند:
\begin{itemize}
  \item \code{backend/.env}
  \item \code{backend/uploads}
  \item \code{backend/storage/logs}
  \item dumpهای SQL، logها و backupها.
\end{itemize}

\subsection{Frontend}
frontend باید با npm build شود:
\begin{codeblock}
cd frontend
npm ci
npm run build
\end{codeblock}

Artifact frontend شامل \code{.next}، \code{public}، \code{package.json} و \code{package-lock.json} است. برای اجرای production:
\begin{codeblock}
npm run start
\end{codeblock}

\subsection{Widget}
در وضعیت فعلی، \code{widget/dist/widget.js} artifact نهایی ویجت است. اگر minify/bundle رسمی اضافه شود، باید خروجی dist در CI تولید و checksum آن ثبت شود. URL نهایی ویجت باید با \code{WIDGET_SCRIPT_URL} هماهنگ باشد.

\section{Docker و Docker Compose}
\subsection{وضعیت فعلی}
در repository فایل \code{Dockerfile} یا \code{docker-compose.yml} وجود ندارد. بنابراین deployment فعلی احتمالا به وب‌سرور سنتی PHP و سرویس Node برای Next.js وابسته است.

\subsection{طرح پیشنهادی Docker}
برای production یا staging قابل تکرار، سه container اصلی پیشنهاد می‌شود:
\begin{itemize}
  \item \code{backend-php}: PHP-FPM یا Apache/PHP برای APIها.
  \item \code{frontend-next}: Node.js برای اجرای Next.js.
  \item \code{mysql}: دیتابیس فقط برای local/staging؛ در production بهتر است سرویس managed یا VM جدا باشد.
\end{itemize}

نمونه ساختار پیشنهادی:
\begin{codeblock}
docker/
  backend/Dockerfile
  frontend/Dockerfile
  nginx/default.conf
docker-compose.yml
\end{codeblock}

\subsection{طرح پیشنهادی Docker Compose}
\begin{codeblock}
services:
  backend:
    build: ./docker/backend
    env_file: ./backend/.env
    volumes:
      - ./backend:/var/www/backend
      - uploads:/var/www/storage/uploads

  frontend:
    build: ./docker/frontend
    environment:
      NEXT_PUBLIC_API_BASE_URL: https://api.example.com/backend/api

  mysql:
    image: mysql:8
    environment:
      MYSQL_DATABASE: ai_chat_saas
      MYSQL_USER: app
      MYSQL_PASSWORD: secret
      MYSQL_ROOT_PASSWORD: root-secret

volumes:
  uploads:
\end{codeblock}

این نمونه در repo فعلی وجود ندارد و باید قبل از استفاده واقعی با نیازهای مسیر فایل، وب‌سرور و SSL تطبیق داده شود.

\section{Kubernetes}
Kubernetes برای MVP ضروری نیست، اما برای SaaS با چند مشتری و نیاز به مقیاس‌پذیری می‌تواند در فاز بعدی استفاده شود. اجزای پیشنهادی:
\begin{itemize}
  \item \code{Deployment} برای frontend و backend.
  \item \code{Service} جدا برای frontend و backend.
  \item \code{Ingress} با TLS و routing دامنه‌ها.
  \item \code{Secret} برای \code{JWT_SECRET}، \code{DB_PASS} و کلیدهای API.
  \item \code{ConfigMap} برای URLها و تنظیمات غیرحساس.
  \item \code{PersistentVolume} یا object storage برای uploadها.
  \item \code{CronJob} برای backup، cleanup و taskهای دوره‌ای.
\end{itemize}

تا زمانی که load و تیم عملیاتی کوچک است، deployment روی VM با Docker Compose ساده‌تر و کم‌ریسک‌تر از Kubernetes است.

\section{CI/CD}
\subsection{وضعیت فعلی}
pipeline رسمی در repository وجود ندارد. بنابراین اجرای build، lint، migration و deploy فعلا باید دستی یا با ابزار بیرونی انجام شود.

\subsection{Pipeline پیشنهادی}
\begin{enumerate}
  \item Checkout کد.
  \item نصب dependencyهای frontend با \code{npm ci}.
  \item اجرای lint frontend با \code{npm run lint}.
  \item اجرای build frontend با \code{npm run build}.
  \item بررسی syntax فایل‌های PHP.
  \item اجرای تست‌های backend/frontend در صورت اضافه شدن.
  \item ساخت artifact backend، frontend و widget.
  \item deploy به staging.
  \item اجرای migration روی staging.
  \item smoke test شامل login، health check دیتابیس، widget config و یک endpoint محافظت‌شده.
  \item approval دستی برای production.
  \item deploy production با نسخه tag شده.
  \item post-deploy health check و مانیتورینگ خطا.
\end{enumerate}

\subsection{نمونه دستورات CI}
\begin{codeblock}
cd frontend
npm ci
npm run lint
npm run build

php -l backend/api/auth/login.php
php -l backend/includes/jwt.php
\end{codeblock}

برای بررسی همه فایل‌های PHP، در CI می‌توان از اسکریپت shell/PowerShell استفاده کرد. این اسکریپت باید روی همه فایل‌های \code{backend/**/*.php} اجرا شود و در صورت خطای syntax pipeline را متوقف کند.

\section{Deployment}
\subsection{Deployment پیشنهادی روی VM}
\begin{enumerate}
  \item دریافت artifact نسخه مشخص از CI یا checkout tag.
  \item نصب dependencyهای frontend با \code{npm ci --omit=dev} یا استفاده از artifact build شده.
  \item کپی backend به مسیر وب‌سرور.
  \item قرار دادن \code{backend/.env} امن خارج از repository یا تزریق از secret manager.
  \item اجرای migrationهای دیتابیس با backup قبل از migration.
  \item تنظیم \code{UPLOAD_STORAGE_ROOT} و permission مناسب.
  \item restart سرویس Node/Next و PHP/Apache/Nginx در صورت نیاز.
  \item اجرای smoke test.
\end{enumerate}

\subsection{مسیرهای پیشنهادی production}
\begin{longtable}{R{42mm} R{88mm}}
\toprule
\textbf{دامنه/مسیر} & \textbf{کاربرد} \\
\midrule
\endhead
\code{https://app.example.com} & frontend پنل Next.js. \\
\code{https://api.example.com/backend/api} & PHP API. \\
\code{https://cdn.example.com/widget.js} & فایل static ویجت. \\
\code{/var/www/ai-chat-saas/backend} & کد backend روی سرور. \\
\code{/var/www/ai-chat-saas/frontend} & کد یا artifact frontend. \\
\code{/var/www/ai-chat-saas-storage} & upload storage خارج از مسیر اجرای کد. \\
\bottomrule
\end{longtable}

\section{Database Migration}
در repo فعلی فقط migration دانش AI دیده شد و migration runner مرکزی وجود ندارد. برای production باید سیاست زیر اجرا شود:
\begin{itemize}
  \item هر تغییر schema در فایل SQL نسخه‌دار داخل \code{backend/database/migrations} ثبت شود.
  \item قبل از اجرای migration در production، backup دیتابیس گرفته شود.
  \item migration ابتدا در staging اجرا شود.
  \item migrationهای destructive باید plan rollback یا اسکریپت جبران داشته باشند.
  \item جدول \code{schema_migrations} برای ثبت migrationهای اجراشده اضافه شود.
\end{itemize}

نمونه اجرای دستی:
\begin{codeblock}
mysql -u app -p ai_chat_saas < backend/database/migrations/2026_07_07_ai_knowledge_tables.sql
\end{codeblock}

\section{Rollback}
\subsection{اصول rollback}
Rollback باید قبل از deploy طراحی شود، نه بعد از شکست. برای هر release باید موارد زیر مشخص باشد:
\begin{itemize}
  \item tag نسخه قبلی سالم.
  \item artifact نسخه قبلی backend، frontend و widget.
  \item وضعیت migration و امکان rollback دیتابیس.
  \item سازگاری نسخه قدیمی frontend با API و دیتابیس جدید.
  \item backup قبل از migration.
\end{itemize}

\subsection{سناریوهای rollback}
\begin{longtable}{R{38mm} R{64mm} R{48mm}}
\toprule
\textbf{سناریو} & \textbf{اقدام} & \textbf{ریسک} \\
\midrule
\endhead
خرابی frontend & بازگرداندن artifact قبلی Next.js و restart سرویس. & پایین، اگر API تغییر نکرده باشد. \\
خرابی backend & بازگرداندن کد PHP قبلی و پاک‌سازی cache وب‌سرور در صورت وجود. & متوسط، وابسته به migration. \\
خرابی widget & بازگرداندن \code{widget.js} قبلی در CDN و کاهش cache TTL. & متوسط، چون روی سایت مشتری cache می‌شود. \\
خرابی migration & restore backup یا اجرای migration جبرانی. & بالا، باید از قبل تست شود. \\
مشکل env/secret & بازگرداندن env قبلی از secret manager و restart سرویس‌ها. & متوسط، نیازمند audit. \\
\bottomrule
\end{longtable}

\section{Monitoring}
\subsection{Health Check}
در حال حاضر \code{backend/api/test-db.php} health check ساده دیتابیس است. برای production بهتر است endpoint اختصاصی \code{/health} اضافه شود که بدون افشای اطلاعات حساس موارد زیر را بررسی کند:
\begin{itemize}
  \item دسترسی PHP runtime.
  \item اتصال دیتابیس با query سبک.
  \item writable بودن storage upload در صورت نیاز.
  \item وضعیت نسخه build.
  \item عدم نمایش stack trace یا جزئیات secret.
\end{itemize}

نمونه پاسخ پیشنهادی:
\begin{codeblock}
{
  "success": true,
  "status": "ok",
  "version": "1.0.0",
  "checks": {
    "database": "ok",
    "storage": "ok"
  }
}
\end{codeblock}

\subsection{شاخص‌های مانیتورینگ}
\begin{itemize}
  \item نرخ خطاهای HTTP 5xx و 4xx.
  \item latency endpointهای login، conversation list، message send و AI reply.
  \item تعداد گفتگوهای ایجادشده در دقیقه.
  \item تعداد رخدادهای rate limit.
  \item خطاهای upload و حجم storage.
  \item مصرف CPU/RAM دیپلوی frontend و backend.
  \item اندازه دیتابیس و سرعت رشد جدول‌های messages، conversations و audit logs.
  \item تعداد خطاهای crawler و AI search.
\end{itemize}

\section{Logging}
\subsection{وضعیت فعلی}
backend از \code{error_log} برای خطاهای امنیتی و runtime استفاده می‌کند. audit مدیریتی نیز در جدول \code{admin_audit_logs} ذخیره می‌شود. در repository، مسیرهای \code{backend/storage/logs} در \code{.gitignore} قرار دارند.

\subsection{پیشنهاد production}
\begin{itemize}
  \item لاگ PHP، وب‌سرور، Node/Next و دیتابیس به log aggregator مرکزی ارسال شود.
  \item logها ساختاریافته و شامل request id، user id، tenant id و endpoint باشند.
  \item secret، password، token و Authorization header هرگز در log ذخیره نشوند.
  \item retention لاگ‌ها بر اساس سیاست تجاری و حریم خصوصی تعریف شود.
  \item alert برای رشد خطای 5xx، شکست login غیرعادی، rate limit شدید و خطاهای upload فعال شود.
\end{itemize}

\section{Tracing}
Tracing توزیع‌شده در repo فعلی پیاده‌سازی نشده است. با توجه به معماری فعلی، ابتدا logging ساختاریافته و request id کافی است. در مرحله بعد می‌توان \en{OpenTelemetry} یا ابزار مشابه را برای ردیابی مسیر درخواست بین frontend، backend و دیتابیس اضافه کرد.

حداقل طراحی پیشنهادی:
\begin{itemize}
  \item تولید \code{X-Request-Id} در reverse proxy اگر client نفرستاده باشد.
  \item عبور دادن request id به PHP و Next.js logs.
  \item ثبت request id در error log و audit context.
  \item نمایش request id در پاسخ خطای production برای پشتیبانی، بدون افشای جزئیات فنی.
\end{itemize}

\section{Auto Deploy}
Auto deploy فعلا در repo پیاده‌سازی نشده است. پیشنهاد:
\begin{itemize}
  \item merge به \code{develop}: deploy خودکار به staging.
  \item tag با الگوی \code{v*}: ساخت release artifact.
  \item deploy production فقط بعد از approval دستی.
  \item smoke test خودکار پس از deploy.
  \item rollback دستی اما سریع با artifact نسخه قبلی.
\end{itemize}

\section{Backup و نگهداری}
\begin{itemize}
  \item backup روزانه دیتابیس و نگهداری چند نسخه زمانی.
  \item backup هماهنگ upload storage با دیتابیس.
  \item تست restore حداقل ماهانه.
  \item پاک‌سازی دوره‌ای logها، فایل‌های موقت و rate limitهای منقضی.
  \item مانیتور فضای دیسک upload و backup.
  \item نگهداری نسخه schema و migrationهای اجراشده.
\end{itemize}

\section{چک‌لیست Release Production}
\begin{enumerate}
  \item branch یا tag release مشخص شده است.
  \item frontend با \code{npm ci} و \code{npm run build} موفق build شده است.
  \item فایل‌های PHP syntax check شده‌اند.
  \item migrationها در staging اجرا و تایید شده‌اند.
  \item backup دیتابیس و upload storage گرفته شده است.
  \item \code{APP_ENV=production} و \code{APP_DEBUG=false} تنظیم شده است.
  \item \code{JWT_SECRET} و \code{DB_PASS} از secret امن تامین شده‌اند.
  \item \code{PANEL_ALLOWED_ORIGINS} و \code{WIDGET_ALLOW_EMPTY_ORIGIN=false} کنترل شده‌اند.
  \item health check و smoke test پس از deploy موفق هستند.
  \item برنامه rollback و artifact نسخه قبلی آماده است.
\end{enumerate}

\section{جمع‌بندی}
DevOps فعلی پروژه برای توسعه و دمو قابل استفاده است، اما برای production نیازمند تکمیل است. اجزای اصلی محصول مشخص و قابل deploy هستند: backend PHP، frontend Next.js، widget static، MySQL و upload storage. با این حال، نبود Docker، نبود pipeline رسمی CI/CD، نبود migration runner، نبود monitoring/tracing مرکزی و نبود health check جامع باید قبل از عرضه گسترده برطرف شود. مسیر پیشنهادی، ابتدا استانداردسازی env و release دستی قابل تکرار، سپس Docker Compose برای staging/production کوچک، و در نهایت CI/CD، monitoring و در صورت نیاز Kubernetes است.

\end{document}
